
%%%%%%%%%%%%%%%%%%%%%%%%%%%%%%%%%%%%%%%%%%%%%%%%%%%%%%%%%%%%
%%%%%%%%%% input macros                           %%%%%%%%%%
%%%%%%%%%%%%%%%%%%%%%%%%%%%%%%%%%%%%%%%%%%%%%%%%%%%%%%%%%%%%



%%%%% normal, italicized greek symbol
%%%%%%%%%%%%%%%%%%%%%%%%%%%%%%%%%%%%%%%%%%%%%%%%%%%%%%%%%%%%

\newcommand{\al}{\alpha}
\newcommand{\be}{\beta}
\newcommand{\ga}{\gamma}
\newcommand{\de}{\delta}
\newcommand{\ep}{\epsilon}
\newcommand{\ze}{\zeta}
%eta
\renewcommand{\th}{\theta}
%iota
\newcommand{\ka}{\kappa}
\newcommand{\la}{\lambda}
%mu
%nu
%xi
%omicron
%pi
%rho
\newcommand{\si}{\sigma}
%tau
\newcommand{\up}{\upsilon}
%phi
%chi
%psi
\newcommand{\om}{\omega}

%Alpha
%Beta
\newcommand{\Ga}{\Gamma}
\newcommand{\De}{\Delta}
%Epsilon
%Zeta
%Eta
\newcommand{\Th}{\Theta}
%Iota
%Kappa
\newcommand{\La}{\Lambda}
%Mu
%Nu
%Xi
%Omicron
%Pi
%Rho
\newcommand{\Si}{\Sigma}
%Tau
%Upsilon
%Phi
%Chi
%Psi
\newcommand{\Om}{\Omega}


%%%%% upright, unitalicized greek symbols
%%%%%%%%%%%%%%%%%%%%%%%%%%%%%%%%%%%%%%%%%%%%%%%%%%%%%%%%%%%%

\newcommand{\ual}{\upalpha}
\newcommand{\ube}{\upbeta}
\newcommand{\uga}{\upgamma}
\newcommand{\ude}{\updelta}
\newcommand{\uep}{\upepsilon}
\newcommand{\uze}{\upzeta}
%eta
\newcommand{\uth}{\uptheta}
%iota
\newcommand{\uka}{\upkappa}
\newcommand{\ula}{\uplambda}
\newcommand{\umu}{\upmu}
%nu
%xi
%omicron
%pi
%rho
\newcommand{\usi}{\upsigma}
%tau
\newcommand{\uup}{\upupsilon}
%phi
%chi
%psi
\newcommand{\uom}{\upomega}

%Alpha
%Beta
\newcommand{\uGa}{\Upgamma}
\newcommand{\uDe}{\Updelta}
%Epsilon
%Zeta
%Eta
\newcommand{\uTh}{\Uptheta}
%Iota
%Kappa
\newcommand{\uLa}{\Uplambda}
%Mu
%Nu
%Xi
%Omicron
%Pi
%Rho
\newcommand{\uSi}{\Upsigma}
%Tau
%Upsilon
%Phi
%Chi
%Psi
\newcommand{\uOm}{\Upomega}




%%%%% macros with greek symbols
%%%%%%%%%%%%%%%%%%%%%%%%%%%%%%%%%%%%%%%%%%%%%%%%%%%%%%%%%%%%


\newcommand{\sisq}{\sigma^{2}}



%%%%%
%%%%%%%%%%%%%%%%%%%%%%%%%%%%%%%%%%%%%%%%%%%%%%%%%%%%%%%%%%%%

%% Short-hand Code
\newcommand{\bs}[1]{\boldsymbol{#1}}
\newcommand{\f}[2]{\frac{#1}{#2}}
\renewcommand{\bar}{\overline}
\newcommand{\floor}[1]{\left\lfloor#1\right\rfloor}
\newcommand{\pa}{\partial}
\newcommand{\mc}[1]{\mathcal{#1}}
\newcommand{\mf}[1]{\mathfrak{#1}}

%% Statistical Stuff
\newcommand{\Indic}[2]{1_{#1}(#2)}
% \newcommand{\E}[1]{\mathbb{E}\left[#1\right]}
\newcommand{\Ewrt}[2]{\mathbb{E}_{#1}\left[#2\right]}
\newcommand{\Var}[1]{\mathbb{V}\left[#1\right]}
\newcommand{\Cov}[2]{\mathrm{Cov}\left[#1,#2\right]}
\newcommand{\CovVec}[1]{\mathrm{Cov}\left[#1\right]}
\newcommand{\Bern}[1]{\mathrm{Bern}(#1)}
\newcommand{\Bin}[2]{\mathrm{Bin}(#1,#2)}
\newcommand{\NegBin}[2]{\mathrm{NegBin}(#1,#2)}
\newcommand{\Geom}[1]{\mathrm{Geom}(#1)}
\newcommand{\Multinom}[1]{\mathrm{Multinom}(#1)} % manually input parameters, e.g. {n, p_{1}, p_{2}, p_{3}}
\newcommand{\Mult}[3]{\mathrm{Multinom}_{#1}(#2,#3)} % same as above, more generic,
\newcommand{\Pois}[1]{\mathrm{Pois}(#1)}
\newcommand{\Unif}[2]{\mathrm{Unif}(#1,#2)}
\newcommand{\Norm}[2]{\mathcal{N}(#1,#2)}
\newcommand{\MVpNorm}[3]{\mathcal{N}_{#1}\left(#2,#3\right)}
\newcommand{\TD}[1]{t_{#1}}
\newcommand{\WDist}[3]{\mathcal{W}_{#1}\left(#2,#3\right)}
\newcommand{\ExpD}[1]{\mathrm{Exp}\left(#1\right)}
\newcommand{\BetaD}[2]{\mathrm{Beta}\left(#1, #2\right)}
\newcommand{\GammaD}[2]{\mathrm{Gamma}\left(#1, #2\right)}
\newcommand{\Weibull}{\mathrm{Weibull}}
\newcommand{\Rayleigh}{\mathrm{Rayleigh}}
\newcommand{\Maxwell}{\mathrm{Maxwell}}
\newcommand{\Logistic}{\mathrm{Logistic}}
\newcommand{\IG}{\mathrm{IG}}
\newcommand{\Gumbel}{\mathrm{Gumbel}}
\newcommand{\Cauchy}{\mathrm{Cauchy}}
\newcommand{\thetavec}{\underline{\theta}}
\newcommand{\etavec}{\underline{\eta}}
\newcommand{\PR}[1]{P\left(#1\right)}
\newcommand{\PRwrt}[2]{P_{#1}\left(#2\right)}
\newcommand{\PRRV}[1]{P\left[#1\right]}
\newcommand{\PRRVwrt}[2]{P_{#1}\left[#2\right]}
\newcommand{\sa}{\sigma\mathrm{-algebra}}
\newcommand{\disju}{\stackrel{\cdot}{\cup}}
\newcommand{\bigdisju}{\stackrel{\cdot}{\bigcup}}
\newcommand{\pisys}{\pi\mathrm{-system}}
\newcommand{\lamsys}{\lambda\mathrm{-system}}
\newcommand{\io}{\mathrm{ i.o. }}
\newcommand{\set}[1]{\left\{#1\right\}}
\newcommand{\oset}[1]{\left[#1\right]}
\newcommand{\tpl}[1]{\left(#1\right)}
\newcommand{\indep}{\perp\!\!\!\perp}
\newcommand{\nindep}{\perp\!\!\!\!\!/\!\!\!\!\!\perp}
\newcommand{\convas}{\stackrel{\mathrm{\scriptsize{a.s.}}}{\longrightarrow}}
\newcommand{\nconvas}{\stackrel{\mathrm{\scriptsize{a.s.}}}{\nrightarrow}}
\newcommand{\convip}[1]{\stackrel{\scriptsize{#1}}{\longrightarrow}}
\newcommand{\nconvip}[1]{\stackrel{\scriptsize{#1}}{\nrightarrow}}
\newcommand{\convlp}[1]{\stackrel{\scriptsize{L_{#1}}}{\longrightarrow}}
\newcommand{\nconvlp}[1]{\stackrel{\scriptsize{L_{#1}}}{\nrightarrow}}
\newcommand{\lpnorm}[2]{\left|\left|#2\right|\right|_{#1}}
\newcommand{\Far}{F^{\leftarrow}}
\newcommand{\convd}{\stackrel{\scriptsize{d}}{\longrightarrow}}
\newcommand{\nconvd}{\stackrel{\scriptsize{d}}{\nrightarrow}}
\newcommand{\eqd}{\stackrel{\scriptsize{d}}{=}}
\newcommand{\logit}{\mathrm{logit}}
\newcommand{\expit}{\mathrm{expit}}
\newcommand{\Func}[1]{{\tt #1}}

%% Mathematical Stuff
\newcommand{\C}{\mathbb{C}}
\newcommand{\D}{\mathbb{D}}
\renewcommand{\R}{\mathbb{R}} % check this
\newcommand{\Q}{\mathbb{Q}}
\newcommand{\N}{\mathbb{N}}
\newcommand{\Z}{\mathbb{Z}}
\renewcommand{\S}{\mathbb{S}}
% \renewcommand{\P}{\mathbb{P}}
\newcommand{\Cinf}{\mathbb{C}_{\infty}}
\newcommand{\CP}{\mathbb{CP}^{1}}
% \newcommand{\RE}{\mathrm{Re}}
% \newcommand{\IM}{\mathrm{Im}}
\newcommand{\re}[1]{\mbox{Re}(#1)}
\newcommand{\im}[1]{\mbox{Im}(#1)}
\newcommand{\du}[1]{\mbox{Du}(#1)}
\newcommand{\Proj}[2]{\mathrm{Proj}_{#1}#2}
\newcommand{\pii}{\pi i}
\newcommand{\prp}{^{\perp}}
\renewcommand{\hat}{\widehat}
\renewcommand{\tilde}{\widetilde}
\newcommand{\st}{^\star}
\newcommand{\epig}[2]{\mbox{epi}_{#1}\left(#2\right)}
\newcommand{\hypog}[2]{\mbox{hyp}_{#1}\left(#2\right)}
\newcommand{\sign}[1]{\mbox{sign}\left(#1\right)}


%%  Type Setting Stuff
\newcommand{\eq}[1]{\begin{equation}#1\end{equation}}
\newcommand{\problem}[1]{\setcounter{equation}{0}\item \textbf{(#1)}}
\newcommand{\sureindent}{$\quad \quad$}
\newcommand{\comment}[1]{\quad \quad (\mathrm{#1})}
\renewcommand{\hbar}{\mathrm{------}}
\renewcommand{\complement}[1]{\bar{#1}}

%% Matrix and vector bolds
\newcommand{\ma}[1]{\textbf{#1}}
\newcommand{\ve}[1]{\textbf{#1}}
\newcommand{\rve}[1]{\boldsymbol{#1}}
% \renewcommand{\null}[1]{\mbox{null}(#1)} % knitr hates this?
\newcommand{\col}[1]{\mbox{range}(#1)}
\newcommand{\norm}[1]{\left|\left|#1\right|\right|}
\newcommand{\specrad}[1]{\rho\left(\mathbf{#1}\right)}
\newcommand{\Cn}{\mathbb{C}^{n}}
\newcommand{\Cnn}{\mathbb{C}^{n \times n}}
\newcommand{\Cmn}{\mathbb{C}^{m \times n}}
\newcommand{\diag}{\mathrm{diag}}
\newcommand{\h}{^{*}}
\renewcommand{\span}[1]{\mathrm{span}\left\{#1\right\}}
\newcommand{\spn}[1]{\overline{\mathrm{sp}}\left\{#1\right\}}
\newcommand{\pspc}{^{\perp}}
\newcommand{\I}{\mathbf{I}}
\newcommand{\tr}{\mathrm{tr} \ }
\renewcommand{\diag}{\mathrm{diag} \ }
\newcommand{\ip}[2]{\left\langle #1, #2\right\rangle}
\newcommand{\twovec}[2]{\left[\begin{matrix}#1 \\ #2 \end{matrix}\right]}
\newcommand{\twomat}[4]{\left[\begin{matrix}#1 & #3 \\ #2 & #4 \end{matrix}\right]}
\newcommand{\threevec}[3]{\left[\begin{matrix}#1 \\ #2 \\ #3\end{matrix}\right]}
\newcommand{\threemat}[9]{\left[\begin{matrix}#1 & #4 & #7 \\ #2 & #5 & #8 \\ #3 & #6 & #9\end{matrix}\right]}
\newcommand{\fourvec}[4]{\left[\begin{matrix}#1 \\ #2 \\ #3 \\ #4\end{matrix}\right]}
\renewcommand{\t}{^{\tau}}
\newcommand{\punc}[1]{#1^{*}}
\newcommand{\curl}[1]{\mbox{curl}(#1)}

%% number theory
\renewcommand{\mod}[2]{#1\,\mathrm{mod}\,#2}
\newcommand{\equivmod}[3]{#1\,=\,#2\,\mathrm{mod}\,#3}
\newcommand{\nequivmod}[3]{#1\,\neq\,#2\,\mathrm{mod}\,#3}
\newcommand{\divides}[2]{#1\,|\,#2}
\newcommand{\ndivides}[2]{#1\,\nmid\,#2}
\renewcommand{\gcd}[2]{(#1,\,#2)}
\newcommand{\divisors}[1]{\mc{D}_{#1}}



%% Other Stuff
\newcommand{\hdefn}[2]{{\color{notsobrightblue}\emph{\href{#1}{#2}}}}
\newcommand\independent{\protect\mathpalette{\protect\independenT}{\perp}}
\def\independenT#1#2{\mathrel{\rlap{$#1#2$}\mkern2mu{#1#2}}}
% \newcommand{\code}[1]{{\tt #1}}


%% colors
\renewcommand{\r}[1]{{\color{red}{#1}}}
\renewcommand{\b}[1]{{\color{blue}{#1}}}
\newcommand{\g}[1]{{\color{green}{#1}}}
\definecolor{BUgreen}{RGB}{24,144,41} % not official BU color
\definecolor{mmGreen}{RGB}{61,153,86}
\definecolor{darkBlue}{RGB}{0,92,141}
\newcommand{\db}[1]{{\color{darkBlue}{#1}}}
\newcommand{\mg}[1]{{\color{mmGreen}{#1}}}

% farver::decode_colour("chocolate")
% \definecolor{firebrick}{RGB}{178,34,34}
% \definecolor{steelblue}{RGB}{70,130,180}
% \definecolor{forestgreen}{RGB}{34,139,34}
% \definecolor{darkorange}{RGB}{139,90,0}  % orange4
% \definecolor{chocolate}{RGB}{75,45,0}  % orange4
% \newcommand{\fb}[1]{{\color{firebrick}{#1}}}
% \renewcommand{\sb}[1]{{\color{steelblue}{#1}}}
% \newcommand{\fg}[1]{{\color{forestgreen}{#1}}}
% \renewcommand{\do}[1]{{\color{chocolate}{#1}}}

% farver::decode_colour("gray50")
\definecolor{grayseventyfive}{RGB}{191,191,191}
\definecolor{grayfifty}{RGB}{127,127,127}
\newcommand{\gsevfive}[1]{{\color{gray}{#1}}}
\newcommand{\gfty}[1]{{\color{grayfifty}{#1}}}


\newcommand{\pl}[1]{{\fontfamily{\sfdefault}\selectfont #1}}
\newcommand{\Rpl}{{\fontfamily{\sfdefault}\selectfont R }}
% \newcommand{\pkg}[1]{\textbf{\pl{#1}}}

\DeclareMathOperator*{\argmin}{\arg\!\min}
\DeclareMathOperator*{\argmax}{\arg\!\max}






